\documentclass[a4paper,oneside,bibliography=totoc]{scrartcl}
% Do not change document class, margins, fonts, etc.

% some useful packages (add more as needed)
\usepackage[utf8]{inputenc}
\usepackage{graphicx}
\usepackage{latexsym}
\usepackage{amsmath}
\usepackage{amssymb}
\usepackage{natbib}
\usepackage{tabularx}
\usepackage{booktabs}
\usepackage{algorithm} % you can modify the algorithm style to your liking
\usepackage{algorithmic}
\usepackage{csquotes}
\renewcommand{\algorithmiccomment}[1]{\hfill\textit{// #1}}
\usepackage[usenames,dvipsnames]{xcolor}
\usepackage[colorlinks,citecolor=Green]{hyperref} % you may change/remove the colors
\usepackage{lipsum} % you do not need this

\bibliographystyle{chicagoa}
\setcitestyle{authoryear,round,semicolon,aysep={},yysep={,}} \let\cite\citep

% example enviroments (add more as needed)
\newtheorem{definition}{Definition} \newtheorem{proposition}{Proposition}

\begin{document}

\subject{Personal Survey} % change to appropriate type
\title{Overview of Knowledge Graphs}
\author{Bagautdin Nukhkadiev}
\date{\today}
\maketitle

\begin{abstract}

Your report must contain an abstract. A good reference for report writing
is~\citet{zobel2014}; we highly recommend that you study this or a similar book
during your studies. He writes the following about the abstract. 

\end{abstract}

\section{Introduction}
\label{ch:intro}

Intro here. 



\section{Literature Review}
\label{sec:related_work}

\subsection{Graph Theory}
% Graphs Definition
Graph Kernels - Graph kernels are a class of methods that let you measure similarity 
between graphs accounting for their structure.

\paragraph{Weisfeiler-Lehman graph kernel algorithm.} 


\subsection{Knowledge Graphs}

Knowledge Graphs are structures made to represent expressive logic and knowledge in a graph. 
Currently Knowledge Graphs are a hype topic that has been on a rise in the last couple of years due to rise in fame of intelligent systems such as LLMs.
Today KGs are used in a range of applications, e.g. enrichment of search results or accurate information injection into LLMs.
To start off, let's walk through the basics of Knowledge Graphs (KGs). 

% Cover the idea of KGs 
The idea of a Knowledge Graph is quite simple: information is structured into triplets called facts, 
e.g. (Albert Einstein, BornIn, German Empire). Elements of these triplets are called (head, relation, tail) or (subject, predicate, object).

% Get into RDF


% OWL 1, OWL 2


% RDFs 

% KG embeddings
\subsection{Knowledge Graph Embeddings}
The notion of embedding information into vector/tensor representations is not new. 
Embeddings play central role in NLP and Information retrieval. 
Thus, embedding Knowledge Graphs seems natural and essential for downstream applications. 
Immediately, we ask ourselves a question:
a) how do we embed the graph structure into a vector
b) how do we ensure that these embedding vectors are able to represent the same level of logic expressivity.
c) how do we measure that. 


KG embedding research worked a lot on building models within different Representation Spaces.
\cite{ji_survey_2022} points out 4 representation spaces: Point-Wise, Complex Vector, Gaussian Distribution and Manifold Space.
\paragraph{Point-Wise Space.} Real number space, basically a go-to usual approach for any algorithm. 
Works like TransE, TransR, HAKE, NTN, DistMult go here

\paragraph{Gaussian Distribution.} Very interesting space, not sure why it is only Gaussian. 
It could be just probabilistic. But the idea that vectors learned are not deterministic but random is great
and I presume along with Point-Wise space should have SOTA models. 


% RDF2vec 
\subsection{RDF2Vec}
As an inspiration from NLP, a Word2Vec embedding method was proposed. There is an idea that h+p=t should hold, thus it makes sense to adapt Word2Vec model proposed by \cite{mikolov_efficient_2013} for KG setting.
This is the motivation for RDF2Vec \cite{groth_rdf2vec_2016}. 

\subsubsection{Graph Walks}
Unfortunately, graphs are not as easy to process as natural language. For that reason a preprocessing step generating sequences is required. 
For that RDF2Vec proposes performing walks on Knowledge Graphs. Original approach suggests Random Walks and Weisfeiler-Lehman Kernels as ways to generate walk sequences. 

\paragraph{Random Walks.}

\paragraph{WL Graph Kernel Sequences.} The walks generated by WL graph kernels are shown in base paper to perform best on AIFB, MUTAG and BGS datasets. 
Such model is called K2V SG because it uses kernel walks and skip-gram model. 

\paragraph{Other Walks.} That point is questionable as to why not use BFS. In this work we shall explore that. 


\subsubsection{Implementations of RDF2Vec}

There are only a couple of implementations of RDF2Vec I've found on the Internet. jRDF2Vec developed by DWS Group and PyRDF2Vec developed by PreDiCT.IDLab. 

\paragraph{jRDF2Vec.} Poorly made package that mixes python and java for no reason. 
Java is used for generating random walks. Then, for some unknown reason a Flask server is created. 
Java application sends the walks file to the server endpoint. 
Server backend executes Word2Vec model from Gensim package.

I find this package to be awfuly designed for 2 main reasons: (i) Redundancy - Flask server, Python and Java are all used for no reason whatsoever. 
(ii) Module does not include essential part of original RDF2Vec paper - Weisfeiler-Lehman graph kernels. 

\paragraph{pyRDF2Vec.} This package is more comprehensive and deserves more attention. 
It has several types of Walks including Weisfeiler-Lehman graph kernels. 
In terms of models it has Gensim's Word2Vec and Fasttext.

Neither of these packages have evaluation or analysis scripts. 



% RDF2vecoa (order-aware)
% Keeps track of where in the path a relation appears
% So it distinguishes:
% direct edge
% vs edge appearing later in a path


\subsection{DLCC}
DLCC is a benchmark dataset for evaluating the performance of RDF2Vec models. 
DLCC introduces a dataset covering 12 test cases that represent a wide range of challenges for RDF2Vec models. 

% Describe the test cases in more detail. 
% Put visualizations of the test cases.
% Use Graphy to visualize the test cases.


\section{Graph Walks}

Walk on graphs are are used for several purposes. 


\section{Graph Sampling}
Sampling for graphs is a technique used to reduce the size of a graph while preserving its properties. 
It is often used in the context of large graphs, where it is not feasible to process the entire graph.







\section{Evaluating RDF2Vec on Complex Datasets}
RDF2Vec models have demonstrated significant potential on benchmark datasets such as DLCC. 
That said, they have not yet been evaluated on more complex datasets.

As a lower baseline we start with a complete nonsense: we initialize RDF2Vec models with random weights and evaluate them on the DLCC dataset.
On the other hand, as a strong baseline to beat we will use TransE-L2 model. 



\begingroup
\newcommand{\D}{\textemdash{}}
\begin{table}[t]
  \centering
  \caption{DLCC (synthetic ontology): binary node classification with a logistic regression on frozen entity vectors. All entries are \textbf{accuracy} on the test set, reported as point estimate $\pm$ bootstrap standard error. Within each setup, columns compare embedding sources.}
  \label{tab:dlcc-rdf2vec}
  \footnotesize
  \setlength{\tabcolsep}{6pt}
  \begin{tabular}{@{}lcccc@{}}
    \toprule
    Test case
      & Random init
      & RDF2Vec
      & ComplEx
      & TransE \\
    \midrule
    TC01 & \D & \D & \D & \D \\
    TC02 & \D & \D & \D & \D \\
    TC03 & \D & \D & \D & \D \\
    TC04 & \D & \D & \D & \D \\
    TC05 & \D & \D & \D & \D \\
    TC06 & \D & \D & \D & \D \\
    TC07 & 0.507 & \D & \D & \D \\
    TC08 & \D & \D & \D & \D \\
    TC09 & \D & \D & \D & \D \\
    TC10 & \D & \D & \D & \D \\
    TC11 & \D & \D & \D & \D \\
    TC12 & \D & \D & \D & \D \\
    \midrule
    Macro avg.
      & \D & \D & \D & \D \\
    \bottomrule
  \end{tabular}
  \vspace{0.4em}
  \par\noindent\small\textit{Note:} Macro avg.\ is the unweighted mean of accuracy over TC01--TC12 for each column. Replace placeholders with $m \pm \sigma$ once runs are complete.
\end{table}

\endgroup



\section{Pretraining a Protograph for RDF2Vec}

We can notice that when training RDF2Vec models we only use the instance graph and not utilize the ontology graph. This is a missed opportunity, as the ontology graph contains valuable information about the structure and semantics of the data. 
To address this, we propose pretraining RDF2Vec models on the ontology graph before fine-tuning them on the instance graph. This approach is inspired by the success of pretraining in other domains, such as natural language processing. 



% 
% What would be the loss function? 
% 




\section{Conclusions}

% \citet{zobel2004} writes:

\blockcquote{zobel2004}{%
  The closing section, or summary, is used to draw together the topics discussed
  in the paper. It should include a concise statement of the paper's important
  results and an explanation of their significance. This is an appropriate place
  to state (or restate) any limitations of the work: shortcomings in the
  experiments, problems that the theory does not address, and so on.

  The conclusions are an appropriate place for a scientist to look beyond the
  current context to other problems that were not addressed, to questions that
  were not answered, to variations that could be explored. They may include
  speculation, such as discussion of possible consequences of the results.

  A \emph{conclusion} is that which concludes, or the end. \emph{Conclusions}
  are the inferences drawn from a collection of information. Write
  "Conclusions", not "Conclusion". If you have no conclusions to draw, write
  "Summary".}

\bibliography{references}


\appendix
\section{Additional Experimental Results}


\section{Proof Details}


\clearpage
\section*{Ehrenwörtliche Erklärung}

Ich versichere, dass ich die beiliegende Bachelor-, Master-, Seminar-, oder
Projektarbeit ohne Hilfe Dritter und ohne Benutzung anderer als der angegebenen
Quellen und in der untenstehenden Tabelle angegebenen Hilfsmittel angefertigt
und die den benutzten Quellen wörtlich oder inhaltlich entnommenen Stellen als
solche kenntlich gemacht habe. Diese Arbeit hat in gleicher oder ähnlicher Form
noch keiner Prüfungsbehörde vorgelegen. Ich bin mir bewusst, dass eine falsche
Erklärung rechtliche Folgen haben wird.

% Declare below which AI tools you used in the process of writing your work,
% including text, image, code, and data generation. If you used a tool for a
% purpose not included in the list yet, add it to the list.
\begin{center}
  \textbf{Declaration of Used AI Tools} \\[.3em]
  \begin{tabularx}{\textwidth}{lXlc}
    \toprule
    Tool & Purpose & Where? & Useful? \\
    \midrule
    ChatGPT & Rephrasing & Throughout & + \\
    DeepL & Translation & Throughout & + \\
    ResearchGPT & Summarization of related work & Sec.~\ref{sec:related_work} & - \\
    Dall-E & Image generation & Figs.~2, 3 & ++ \\
    GPT-4 & Code generation & functions.py & + \\
    ChatGPT & Related work hallucination & Most of bibliography & ++ \\
    \bottomrule
  \end{tabularx}
\end{center}

\vspace{2cm}
\noindent Unterschrift\\
\noindent Mannheim, den XX.~XXXX 2024 \hfill

\end{document}
